\documentclass{newlayout}
%Bitte hier den enstprechenden Ort einsetzen z.B. Braunschweig und die Akademienummer
\Akademie{Braunschweig}{2015}{1}

\usepackage[english, ngerman]{babel}
\usepackage{misc}
\usepackage{multicol}

%\usepackage{amsmath}%wird automatisch durch newlayout.cls geladen
\usepackage{amsfonts}

\usepackage{blindtext}
\usepackage{url}
\def\UrlBreaks{\do\a\do\b\do\c\do\d\do\e\do\f\do\g\do\h\do\i\do\j\do\k\do\l%
\do\m\do\n\do\o\do\p\do\q\do\r\do\s\do\t\do\u\do\v\do\w\do\x\do\y\do\z\do\0%
\do\1\do\2\do\3\do\4\do\5\do\6\do\7\do\8\do\9\do\-\do\_\do\/\do\%}
\urlstyle{same}

% hinzugefügt, um Fehler 'pdfTeX error (font expansion): auto expansion is only possible with scalable' zu vermeiden
\usepackage{lmodern}
\setkomafont{descriptionlabel}{\normalfont\bfseries}
\addtokomafont{paragraph}{\normalfont}
\usepackage{footnote}
\usepackage[flushmargin,hang,ragged]{footmisc}
\deffootnote{1em}{1em}{%
\textsuperscript{\thefootnotemark\ }
}
%\setlength{\abovedisplayskip}{5pt}
%\setlength{\belowdisplayskip}{5pt}

\newcommand{\media}{media/}

%%%%%Mathe-Definitionen
\newtheorem{Def}{Definition}
\newtheorem{Sat}{Satz}
\newtheorem{Bew}{Beweis}
\newtheorem{Thm}{Theorem}

\setlength\abovedisplayshortskip{0pt}
\setlength\belowdisplayshortskip{0pt}
\setlength\abovedisplayskip{3pt}
\setlength\belowdisplayskip{3pt}
%%%%Ende Mathe-Definitionen

\begin{document}
%\input{titel}
\setcounter{page}{3}

\setcounter{tocdepth}{1}
\tableofcontents

\setcounter{secnumdepth}{1}

\setcounter{page}{7}
\setcounter{chapter}{0}

%Angabe, bis zu welcher Stufe die sections im Text nummeriert werden sollen.
\settocdepth{2}

\course{1.2}{Learning from Data}%%% 
\begin{coursetitle}
  \centerline{Learning from Data} 
  \bigskip
  \Large \centerline{Mathematische Theorien und Methoden für das digitale Zeitalter}
  \bigskip
 \includegraphics[width=.9\columnwidth]{media/Titelbild-fehlt.png}
 \label{fig:meinbild}
  \bigskip
\end{coursetitle}
%\begin{dsafigure}
%\begin{center}
%\includegraphics[width=.9\columnwidth]{Titelbild-fehlt.png}
%\caption{meine Bildunterschrift}
%label{fig:meinbild}
%\end{center}
%\end{dsafigure}
%\documentclass{article}

\usepackage{ dsfont }
\usepackage{amsmath}
\usepackage[utf8]{inputenc}
%\usepackage{mathtools}

\begin{document}

\section{Analysis}


Analysis ist ein wichtiges Teilgebiet der Mathematik. Sie ist Grundlage für Optimierung und somit auch von großer Bedeutung beim Maschinellen Lernen. \\
Im Kurs war f\"ur uns besonders die Stetigkeit und Differenzierbarkeit von ein- und mehrdimensionalen Funktionen wichtig. \\
Um uns Klarheit \"uber diesen Bereich der Mathematik zu verschaffen, l\"osten wir - unter anderem - die folgenden Aufgaben.

\vspace{15pt}

\textbf{Aufgabe 1:} \\
Aufgabenstellung:
Beweise, dass eine Funktion $f:\mathds{R} \rightarrow \mathds{R}$, welche differenzierbar ist an der Stelle $x \in R$ ist auch stetig an bei $x \in R$ ist.\\ %Satz ist noch nicht optimal formuliert

Lösung:

Zuerst überlegen wir, welche Bedingungen bereits in der Aufgabenstellung gegeben sind. \\
Eine Funktion ist an der Stelle x differenzierbar, wenn gilt:\\ \\

$f(x+h)=f(x)+l_x(h)+ r(h)$ mit $\lim\limits_{h \rightarrow 0} \frac{r(h)}{h}= 0$ \\ \\

Diese Funktion wäre außerdem stetig, wenn $\lim\limits_{w \rightarrow x} f(w)=f(x)$.\\
Da die Funktion $f(x+h)$ äquivalent zur Funktion $f(w)$ sein soll, ersetzen wir in der Bedingung für die Stetigkeit $f(w)$ mit der gesamten Funktion von $f(x+h)$.  
Somit gilt: \\ \\
$\lim\limits_{h \rightarrow 0} f(x)+l_x(h)+r(h)=f(x)$. \\ \\

Damit die Bedingung der Stetigkeit erfüllt ist, muss der Grenzwert dieser Funktion gleich dem Funktionswert sein. Um diese Bedingung zu erfüllen, müssen wir beweisen, dass die Teile der Funktion $l_x(h)$ und $r(h)$ gegen Null gehen. Wir beginnen mit dem Teil $r(h)$.  Würde $r(h)$ nicht gegen Null gehen, so würde $\frac{r(h)}{h}$ nicht gegen Null gehen. Die Funktion $r(h)$ muss gegen Null gehen, da die Bedingung der Differenzierbarkeit gilt. Jetzt müssen wir zeigen, dass $l_x(h)$ auch gegen Null geht. Dieser Teil der Funktion kann auch als $f'(x) \cdot h$ dargestellt werden. Da wir annehmen, dass h gegen Null geht und somit ein Faktor von $l_x(h)$ Null ist, wird die Funktion Null. Damit ist bewiesen, dass $\lim\limits_{h \rightarrow 0} f(x)=f(x)$. \\

\vspace{15pt}

\textbf{Aufgabe 3:}\\
Bei der nachfolgenden Aufgabe soll bewiesen werden, dass es immer möglich ist für eine differenzierbare Funktion f einen Skalar $\lambda$ zu finden, damit folgende Aussage gilt: \\ \\
$(1) f(x + \lambda \nabla f(x)) \leq f(x)$ \\ \\
Manche von Ihnen werden sich jetzt fragen, was das Zeichen $\nabla$ bedeutet. Es steht für den Gradienten, welcher im Grunde die Ableitung einer Funktion darstellt.
Dadurch lässt sich die Ungleichung (1) in folgende Form umwandeln: \\ \\
$f(x) + Df(x)h + r(h) \leq f(x)$  \\ \\
Hierbei ist h = $\lambda \nabla$ f(x) und r(h) das Restglied %Muss eventuell erklärt werden 
. \\
Das Ganze kann man noch weiter umschreiben, indem man das Skalarprodukt bildet: \\ \\
$f(x) + \lambda <\nabla f(x), \nabla f(x)> + r(h) \leq f(x)$ \\ \\
Anschließend kann man von der Regel Gebrauch machen, dass sich ein Skalarprodukt, das in der Form $<u,u>$ vorliegt, in die Auspr\'agung der Euklidischen Norm von $||u||^2$ umgewandelt werden kann. \\
Daraus folgt: \\ \\
$f(x) + \lambda \cdot || \nabla f(x)||^2 + r(h) \leq f(x)$ \\ \\
Um nun die obige Ungleichung (1) zu beweisen, muss man zeigen, dass \\ $\lambda ||\nabla f(x)||^2 + r(h) \leq 0$ ist.\\
Dafür wird $\lambda$ ausgeklammert. Anschließend erhält man:\\
Da $||\nabla f(x)||$ quadriert wird, wird dieser Term automatisch positiv. Da man den gesamten Term aber negativ haben will, wählt man $\lambda < 0.$ Schließlich kennt jeder die Regel, dass ein Produkt negativ wird, falls eine ungerade Zahl von Faktoren negativ ist. Da wir $\lambda$ negativ gewählt haben, muss \\ $||\nabla f(x)||^2$$(2) \lambda (||\nabla f(x)||^2 + \frac{r(h)}{\lambda}$ positiv sein.\\
Bei $||\nabla f(x)||^2||$ stellt dies aus den eben genannten Gr\"unden kein Problem dar. Deswegen muss nur noch gezeigt werden, dass es mindestens einen Wert von $r(h)$ gibt, der gr\"oßer oder gleich 0 ist. Ansonsten würde der Wert in der Klammer negativ werden k\"onnen und damit die unspr\"ungliche Aussage (1) widerlegen.\\
Damit wir zeigen k\"onnen, dass $r(h) >= 0$ ist, formten wir zuerst den Term (2) um:\\

$\lambda (||\nabla f(x)||^2 + \frac{r(h)}{\lambda})$ \\

$\Leftrightarrow \lambda(||\nabla f(x)||^2 + \frac{r(\lambda \nabla f(x))}{\lambda})$ \\

Bei dieser Umformung wurde h wieder mit $\lambda\nabla f(x)$ ersetzt.

$\leftrightarrow \lambda ||\nabla f(x)|| \cdot (||\nabla f(x)|| + \frac{r(\lambda \nabla f(x))}{\lambda || f(x)||})$


Durch diese Umformungen schlussfolgerten wir, dass f\"ur die Funktion $h(x)$ folgendes gelten muss:\\
(3) $h(\lambda)=||\nabla f(x)|| + \frac{r(\lambda\nabla f(x))}{\lambda||\nabla f(x)||} >=0$ \\

Da in $||\nabla f(x)$ kein $\lambda$ steht ist dieser Term konstant. Gleichzeitig wird durch das Bilden der euklidischen Norm festgelegt, dass ein Wert größe/gleich null vorliegt. \\ %der Satz ist noch nicht wiklich schön formuliert
Deswegen muss nur noch gezeigt werden, dass dder zweite Summand ebenfalls positiv werden kann. Dazu bildeten wir den Grenzwert. \\ %Die Zeitformen stimmen sehr oft nicht
$\lim\limits_{h \rightarrow 0}{\frac{r(\lambda \nabla f(x))}{\lambda ||\nabla f(x)||}} = \frac{r(\lambda \nabla f(x))}{||\lambda  \nabla f(x)||}$ \\ %hier kommt noch kein h vor. Umschreiben?

Man kann $\lambda \nabla f(x)$ wieder durch h ersetzen, woraus folgt: \\
$\lim\limits_{h \rightarrow 0} \frac{r(h)}{h}$ \\

Der Grenzwert des Restbetrages strebt gegen null. Daraus lässt sich schließen, dass es einen Punkt gibt, ab welchem die Ungleichung der Funktion $h(\lambda)$ (3) erfüllt ist und damit auch die anfängliche Aussage (1). \\ %ist hier h(\lambda) und (3) beies nötig?
\vspace{15pt}

\textbf{Aufgabe 6:} \\

In der n\"achsten Aufgabe soll gezeigt werden, dass zwei differenzierbare Funktionen, die vom $\mathds{R}^m$ in den $\mathds{R}$ abgebildet werden, mit $h(x) = f(x) \cdot g(x)$ die Produktregel bedienen. Diese lautet: $Dh(x) = f(x) \cdot Dg(x) + Df(x) \cdot g(x)$.


\end{document}



%%\documentclass{article}

%\usepackage[ngerman]{babel}
%\usepackage[utf8]{inputenc}
%\usepackage{color}
%\usepackage[a4paper,lmargin={2cm},rmargin={2cm},
%tmargin={2.5cm},bmargin = {2.5cm}]{geometry}
%\usepackage{amssymb}
%\usepackage{amsmath}
%\usepackage{graphicx}
%\usepackage{multicol}
%\usepackage{amsthm}

%\setlength{\columnsep}{1cm}
%\newtheorem{defi}{Definition}
%\newtheorem*{defi*}{Definition}

%\begin{document}
%\begin{multicols*}{2}

\section{Vektoren}
In der linearen Algebra nehmen Vektoren eine zentrale Rolle ein. Ein Vektor wird in der Schule als eine Menge an Pfeilen gelehrt, die parallel, gleichgerichtet und gleich lang sind. Betrachten wir jedoch nicht nur den dreidimensionalen Raum, sondern $\mathbb{R}\textsuperscript{n}$, so ist die Vorstellung eines Pfeils nicht immer m\"oglich . Dies bedingt die Notwendigkeit einer anderen Definition.
\newline
\indent Ein Vektor ist eine Aufz\"ahlung von Objekten und beschreibt eine Verschiebung. Eine solche Aufz\"ahlung entspricht der Definition eines Tupels. Entscheidend bei Tupeln ist die Reihenfolge der Objekte, die auch mehrfach vorkommen k\"onnen.
\begin{equation*}
\text{n-Tupel } (a_1,a_2,\dots,a_n)
\end{equation*}
\subsection{Vektor- \& Unterraum}
Vektoren bilden die Elemente eines Vektorraumes V. Addieren wir zwei Vektoren eines Vektorraumes oder multiplizieren wir sie mit einem Skalar, so ist die Summe bzw. das Produkt ebenfalls ein Element des Vektorraumes.
\begin{align*}
(1)\text{ }u, v \in V,\text{ } u+v \in V  \\
(2)\text{ }u \in V, \lambda \in \mathbb{R} ,\text{ }\lambda u \in V
\end{align*}
Ein Unterraum ist eine Teilmenge eines Vektorraumes. Es gelten f\"ur sie die oben genannten Eigenschaften eines Vektorraumes.
\begin{align*}
U\subseteq V, \text{wenn gilt:}	\\
(1)\text{ }u, v \in U,\text{ } u+v \in U  \\
(2)\text{ }u \in U, \lambda \in \mathbb{R} ,\text{ }\lambda u \in U
\end{align*}
\subsection{Span}
Alle m\"oglichen Vektoren eines Vektorraumes werden durch den sogenannten Span dargestellt. Der Span ist die Menge aller m\"oglichen Linearkombinationen der Basisvektoren.
\begin{align*}
\text{Span }(V_1,\dots,V_k) = \{U \in V : U = \lambda_1v_1+\dots+\lambda_kv_k\} \\
\text{ mit } \lambda_1,\dots,\lambda_k \in \mathbb{R}
\end{align*}
\subsection{Basis \& Dimension}
Die Menge der Basisvektoren wird Basis eines Vektorraumes genannt. Jeder Vektorraum besitzt eine Dimension p, die durch die Anzahl der Basisvektoren bestimmt wird.
\subsubsection*{\"Ubung}
Beweise, dass ein Span immer ein Vektorraum ist.
\begin{proof}
Wir nehmen an: 
\begin{align*}
U \in Span(V_1,\dots,V_k)\\
V \in Span(V_1,\dots,V_k)\\
\\
U = \lambda_1V_1+\dots+\lambda_kV_k\\
V = \alpha_1V_1+\dots+\alpha_kV_k\\
\end{align*}
Additionsregelung bei Vektorr\"aumen (s.(1))
\begin{align*}
	&W_1 = U+V\\
	&=\lambda_1V_1+\dots+\lambda_kV_k+\alpha_1V_1+\dots+\alpha_kV_k\\
	&= (\lambda_1\alpha_1)V_1+\dots+(\lambda_k\alpha_k)V_k
\end{align*}
Multiplikationsregelung bei Vektorr\"aumen (s.(2))
\begin{align*}
	&W_2 = \lambda_1V_1+\dots+\lambda_kV_k\\
	&\lambda W_2= (\lambda \lambda_1)V_1+\dots+(\lambda\lambda_k)V_k
\end{align*}
\end{proof}
\subsection{Lineare Unabh\"angigkeit}
$v_1,\dots,v_k$ mit $v_i \in V$ sind linear unabh\"angig, falls $\lambda_iv_i+\dots+\lambda_kv_k = 0$ nur f\"ur $\lambda_i=\dots=\lambda_k= 0$ gilt. 

%\end{multicols*}
%\end{document}

%\subsection{Grundlagen der Optimierung}
Bei der Optimierung betrachtet man vor allem Funktionen, welche Elemente eines $n$-dimensionalen Vektorraumes auf reelle Zahlen abbilden:

\begin{equation*}
f: \mathbb{R}^n \rightarrow \mathbb{R}
\end{equation*}

Beispielhafte Anwendungszwecke wären die Minimierung oder Maximierung der Kosten bzw. der Einnahmen eines Unternehmens. Diese würden dem reellen Funktionswert der Funktion entsprechen und ergäben sich aus den Eigenschaften des Unternehmens wie Gehalt, Anzahl der Mitarbeiter, Marketing-Ausgaben usw., welche man zusammenfassend als Vektor eines $n$-dimensionalen Vektorraumes darstellen könnte.

\subsection{Regression}

\begin{Thm}[Methode der kleinsten Quadrate]
Beispiel für die Lösung eines solchen Minimierungsproblems wäre die Methode der kleinsten Quadrate. Gegeben ist hierbei eine lineare Funktion $f$ und Punkte $P_i=(x_i,y_i)$, welche näherungsweise auf der Geraden liegen:

\begin{align*}
f(x) &= ax \\
y_i &= f(x_i) + \epsilon_{i}
\end{align*}

\begin{dsafigure}
\begin{center}
\includegraphics[width=0.5\textwidth]{\media Grafik-Optimierung_LineareRegression.pdf}
\caption{Konvexe Funktion}
\label{figure:Grafik-Optimierung_LineareRegression.pdf}
\end{center}
\end{dsafigure}

Wenn nun der Abstand $\epsilon_i$ minimiert werden soll, dann gilt:

\begin{equation*}
\min h(x)=\frac{\sum_{i=1}^n(y_i-ax_i)^2)}{e(y_i,x_i)}
\end{equation*}

Hierbei könnte $e^2,\mid e\mid, h(e)$ exemplarische Möglichkeiten seien. Zur Berechnung dieses Minimums gibt es nun verschiedene Möglichkeiten:

\begin{enumerate}
\item Gewöhnliche Tiefpunktberechnung \\
z.~B.$\qquad h'(a)=0\qquad \wedge\qquad h''(a)>0$
\item Gradientenabstiegsmethode \\
$a^{t+1}  =a^{t} - \lambda \cdot \delta f(a^{t})$ \\
Man findet also ein $\lambda > 0$, sodass
$f(a^{t+1}) < f(a^{t})  \qquad \mid \delta f(a^{(t)}) \neq 0$ \\
$g^{+}(\lambda)=f(a^+-\lambda\Delta f(a^t))$ \\
$\lambda^+=argmin(g^+(\lambda))$
\end{enumerate}

\end{Thm}

\begin{Def}[Lokales und globales Minimum]
Für $f: D\rightarrow \mathbb{R}$ ist $x\in D$ ein lokales Minimum, wenn mindestens eine Umgebung $N$ existiert, sodass \\$\forall_{y\in N}$ gilt: $f(y)>=f(x)$ \\
$x\in D$ ist dann ein globales Minimum, falls $N=D$.
\end{Def}

\subsection{Konvexität}

\begin{Def}[Konvexe Funktion]
Eine Funktion $f: \mathbb{R}^n\rightarrow\mathbb{R}$ ist dann eine konvexe Funktion, falls $\forall_{x,y\in D}$ gilt: 
$f(\lambda x+ (1 - \lambda)y) \leq \lambda f(x)+(a-\lambda)f(y) \qquad \mid \forall_{y\in[0,1]}$

\begin{dsafigure}
\begin{center}
\includegraphics[width=0.45\textwidth]{\media Grafik-Optimierung_KonvexeFunktion.jpg}
\caption{konvexe Funktion}
\end{center}
\end{dsafigure}

\end{Def}

\subsubsection{Beispiel zur Konvexität}

\begin{align*}
f(u) &=\mid\mid u\mid\mid_2^2\qquad\mid u\in\mathbb{R}\\
f(\lambda u+(1-\lambda)v) &=\mid\mid\lambda u+(1-\lambda)v\mid\mid_2
\leq\mid\mid\lambda u\mid\mid_2+\mid\mid(1-\lambda)v\mid\mid_2
\leq\lambda\mid\mid u\mid\mid_2+(1-\lambda)\mid\mid v\mid\mid_2
\end{align*}
%\subsection{Konvexe Mengen}

\paragraph{Einführung}

Im Rahmen der Optimierung von konvexen Funktionen ist es erforderlich, den Begriff der konvexen Menge einzuführen. Um diese Thematik anschaulich darzustellen, verwenden wir zunächst verschiedene geometrische Figuren in Abbildung \ref{figure:Grafik-Optimierung_KonvexeMengen}, von denen einige konvex sowie andere wiederum nicht konvex sind. \\
Ist es eine direkte Verbindungsstrecke zwischen zwei beliebigen zu finden, die selbst ebenfalls in der Menge liegt, so ist die Menge konvex.

\begin{dsafigure}
\begin{center}
\includegraphics[width=0.4\textwidth]{\media Grafik-Optimierung_KonvexeMengen.pdf}
\label{figure:Grafik-Optimierung_KonvexeMengen}
\caption{Beispiele konvexer Mengen}
\end{center}
\end{dsafigure}

\begin{Def}[Konvexe Menge]

Eine Menge $X$ heißt konvex, falls für alle $x, y \in X$ und $\lambda \in \mathbb{R}$ und $\lambda \in [0,1]$ gilt:

\begin{equation}
\lambda x + (1 - \lambda) \in X
\end{equation}

\end{Def}

\subsection{Beispiele konvexer Mengen}

\paragraph{Kugel}

Ein Beispiel für eine konvexe Menge ist die Menge $S$ der Vektoren, die eine Kugel mit dem Radius $r = 1$, deren Mittelpunkt im Ursprung liegt, beschreibt:

\begin{equation*}
S = \{x \in \mathbb{R}^{n} | \lVert x \rVert \le 1\}
\end{equation*}

\paragraph{Quadrant}

Auch die Quadranten des kartesischen Koordinatensystems lassen sich durch eine konvexe Menge erfassen. So gilt z.~B. für den ersten Quadranten:

\begin{equation*}
L = \{ x \in \mathbb{R}^{n} | x_1 \ge 0, x_2 \ge 0, ..., x_n \ge 0\}
\end{equation*}

\paragraph{Eistüte}

Bei einer Eistüte bzw. einem quadratischen Kegel handelt es sich um ein weiteres Beispiel für eine konvexe Menge, da sie folgendermaßen für $x \in \mathbb{R}^{n-1}$ und $t \in \mathbb{R}$ beschrieben werden kann:

\begin{equation*}
Q = \{(x, t) \in \mathbb{R}^{n} | \lVert x\rVert ^{2} \le t\}
\end{equation*}

\paragraph{Box}

Die Menge von Vektoren in $B \subseteq \mathbb{R}$ bilden eine konvexe Menge in Form einer Box mit der Seitenlänge 2, deren Mittelpunkt im Ursprung liegt:

\begin{equation*}
B = \{x \in \mathbb{R}^{n} | |x_{i}| \le 1, i = 1, ..., n\}
\end{equation*}

\subsection{Konvexe Programme}

\paragraph{Einführung}

Bei einem konvexen Programm handelt es sich um die Optimierung einer Funktion $f$, die vom Parameter $x$ abhängig ist. Wir versuchen diese Funktion unter Berücksichtigung von Nebenbedingungen zu minimieren. Hierbei sind sowohl die Zielfunktion als auch die Menge der Punkte, die die Nebenbedingungen erfüllt, konvex.

\begin{Def}[Konvexes Programm]
Sei $f: \mathbb{R}^{n} \rightarrow \mathbb{R}$ eine konvexe Funktion und $X \subseteq \mathbb{R}^{n}$, so ist das konvexe Optimierungsproblem bzw. Programm die Minimierung der Funktion $f(x)$ unter der Nebenbedingung, sodass $x \in X$.
\end{Def}

\paragraph{Lineare Programme}

Lineare Programme lassen sich durch lineare Funktionen beschreiben, z.~B. $f(x) = \langle c, x \rangle$. Darüber hinaus sind eine lineare Abbildungsmatrix $A$, z.~B. $A = \begin{pmatrix}1 & 0 \\ 0 & 1 \end{pmatrix}$, sowie ein Vektor $b$ und ein weiterer Vektor $c$, z.~B. $c = \begin{pmatrix}1 \\ 0 \end{pmatrix}$, gegeben. $x$ ist dann zulässig, wenn $x \in X$ mit $X: Ax \le b$ bzw. $X = \{x \in \mathbb{R}^{n} | Ax \le b\}$ erfüllt ist. \\
Ein Beispiel für ein lineares Programm wäre:

\begin{equation*}
f(x) = \langle \begin{pmatrix}1 \\ 0 \end{pmatrix}, \begin{pmatrix}x_{1} \\ x_{2} \end{pmatrix} \rangle = x_{1}
\end{equation*}

Für die Nebenbedingungen gilt:

\begin{equation*}
\begin{pmatrix}1 & 0 \\ 0 & 1 \end{pmatrix}
\end{equation*}

Für den Wertebereich des linearen Programms folgt: $x_{1} \le 1$, $x_{2} \le 1$.

\paragraph{Quadratische Programme}

Im Vergleich zu einem linearen Programm weist z.~B. eine Kostenfunktion im quadratischen Programm eine quadratische Form auf. \\
Beispielsweise sei ein quadratisches Programm $f(x) = x^{T}Qx$ mit $Q = \begin{pmatrix}1 & 0 \\ 0 & 1 \end{pmatrix}$. $f$ gilt es nun unter der Bedingung $x \in X$ für $X: Ax \le b$ zu minmieren:

\begin{align*}
f(x) &= x^{T}Qx \\
&= \begin{pmatrix}x_{1}, x_{2}\end{pmatrix} \begin{pmatrix}1 & 0 \\ 0 & 1\end{pmatrix} \begin{pmatrix}x_{1} \\ x_{2}\end{pmatrix} \\
&= \begin{pmatrix}x_{1}, x_{2}\end{pmatrix} \begin{pmatrix}x_{1} \\ x_{2}\end{pmatrix} \\
&= x_{1}^{2} + x_{2}^{2}
\end{align*}
%\subsection{Gradientenmethode}

Folgender Algorithmus implementiert die Gradientenmethode, die einen wichtigen Teil des Gebietes der Optimierung darstellt. Funktionieren tut sie wie folgt:
\begin{algorithmic}[1]

\Procedure{$\mathbf{calculateMinimum}$}{}
   \\$F(x) = (x^2-2)^2$
\\F'(x) = 4x$^3$-8x

\\
\\x = 10.0
\\$\lambda = 0.001$

\For{$i = 1, \dots, m$}
 \State $\lambda = \lambda+0.001\cdot i$
  \For{$j = 1, \dots, n$}
    \State $x = x - \lambda \cdot F'(x)$
  \EndFor
\EndFor
 
 
\\Print x\EndProcedure
\Statex
\end{algorithmic}

\paragraph{Erklärung}
$F(x)$ beschreibt die Funktion, deren globales Minimum wir finden wollen. $F'(x)$ ist dementsprechend die Ableitung der Funktion $F(x)$. Mit $x=10.0$ setzen wir den Schätzwert, ab dem optimiert wird. $\lambda$ bekommt einen niedrigen Wert deklariert, damit man sich in kleinen Schritten dem Minimum annähern kann. Dies wiederum geschieht in zwei For-Schleifen, deren Inhalt in diesem Fall je 8 mal durchlaufen wird. In der äußeren Schleife sorgen wir dafür, dass unser Lambda größere Werte annimmt. In der inneren For-Schleife wird die Schätz-Variable $x$ mit Hilfe der Ableitung und $\lambda$ angepasst. Daraus folgt:
\begin{equation*}
 f(x+\lambda \nabla f(x)) \leq f(x)
 \end{equation*}



%\newcommand{\media}{../../media/}


\section{Machine Learning}
\author{Farhadiba Mohammed, Dennis Kempf, David Steinmann}
Machine Learning ist ein Konzept, bei dem der Computer durch Algorythmen in der Lage ist Probleme selbstständig zu lösen. Dabei erkennt der Computer nach einer gewissen Lernphase Gesetzmässigkeiten, durch die er dann in der Lage ist, neue, ihm unbekannte Datensätze bezüglich eines bestimmten Problems zu lösen.

Dass der Computer die Gesetzmässigkeiten erkennen kann, werden verschiedene Systeme genutzt, doch zunächst einmal sollten wir uns mit den Begriffen Klassifikation und Regression vertraut machen.


\subsection{Klassifikation und Regression}
\author{David Steinmann}
Bei Klassifikation und Regression handelt es sich um zwei verscheidene Arten von Problemen, beziehungsweise Möglichkeiten zur Problemlösung im Bereich Machine Learning.

\subsubsection{Klassifikation}
Bei der Klassifikation sorgt der Algorythmus am Ende dafür, dass der Datensatz klassifiziert wird. Das heißt, der Datensatz wird vom Computer in verschiedene Klassen eingeteilt. Die Klassen müssen vorher vom Mensch entschieden werden. Der Computer klassifiziert die Datensätze dann anhand der Attribute des entsprechenden Datensatzes. Die klassifizierten Datensätze enthalten dann meistens eine weitere Dimension, in der die Klassifizierung anhand einer Zahl gespeichert ist.

\subsubsection{Regression}
Bei der Regression geht es darum, den entsprechenden Datensätzen am Ende einen festen Wert zuzuordnen. Der Unterschied zur Klassifikation besteht darin, dass der zugeordnete Wert nicht für eine Klasse steht, sondern ein konkretes Ergebnis beinhaltet. Ein Beispiel wäre der Preis von einer Wohnung. Diese wurde mit vielen verscheidenen Attributen in die Datenbank eingespeichert (z. B. Größe, Lage ...) der Endwert wäre dann zum Beispiel der konkrete Preis und nicht eine Klassifizierung in "`teuer"', "`mittelteuer"', "`billig"'. 
%Noch die Mathematische Regression (methode der kleinsten Quadrate)

\subsection{Support Vector Machines}
\author{David Steinmann}
Support Vector Machines oder kurz SVMs sind ein Konzept zum Lösen von Machine Learning Problemen.
Dabei werden für alle Datensätze zürst gewisse Features festgelegt. Diese Features entsprechen den Attributen des Datensatzes. Je nach Problem kann es unterschiedlich viele Features geben, doch allgemein kann man sagen, dass, je mehr Features es sind die Genauigkeit des Programms genaür wird.

Diese Features werden normalerweise in einem Vektor für den entsprechenden Datensatz gespeichert.
Dieser Vektor wird $x_{i}$ genannt wobei $i = 1, ..., n$ den Datensatz beschreibt mit  $x \in  \mathbb{R}^m$.
Zusätzlich hat jeder Datensatz noch einen weiteren Wert, welcher das >>Ergebnis<< des Datensatzes beschreibt, der $y_{i} \in \mathbb{R}$ oder auch >>Label<< genannt wird.

Die Features und das Label zusammen beschreiben einen Punkt im n-dimensionalen Koordinatensystem, wobei $n = m + 1$, da zu den m Dimensionen von x noch die eine Dimension von y dazukommt.

Durch die Features ist es dann möglich den Datensätzen einen bestimmten Ort im n-dimensionalen Koordinatensystem zuzuordnen. Diese können dann auf Grund ihrer räumlichen Anordnung klassifiziert werden. Dabei soll der Abstand von beiden Gruppen von Punkten zu dem Trennobjekt maximal sein. Dadurch kann die Klassifizierung optimal durchgeführt werden. Andernfalls würden schon geringe Abweichungen von den Test-Datensätzen reichen, das ein Datensatz falsch klassifiziert würde. (\ref{SVM2}) 


\begin{dsafigure}
\begin{center}
	\includegraphics[width=0.5\textwidth]{\media svm.jpg}
	\caption{Da der Abstand zwischen den unterschiedlich klassifizierten Datensätzen maximiert werden soll, gilt die rote und nicht die blaue Linie als Trennelement.}
	\label{SVM1}
	\end{center}
\end{dsafigure}


Um die Trennlinie zu optimieren, gibt es auch eine mathematische Beschreibung. Da sich Ebenen auch in der Form $\langle x, w \rangle = b $ beschreiben lassen, folgt, dass wir die Datensätze auch etwas anders betrachten können (\ref{SVM2}).
Das daraus folgende mathematische Programm sieht wie folgt aus: 

\begin{align*}
	\operatorname{maximiere}_{w,b} \quad & \frac{2}{\norm{(w)}} \\
	\operatorname{sodass} \quad &y_{i}(\langle w,x_{i} \rangle -b) \geq 1 \quad \forall  i = 1,...,n ; \\
	&x \in \mathbb{R}^n; \quad	y \in \{1;-1\}
	 %\label{SVMProblemMax}
\end{align*}

Unter der Vorraussetzung, dass alle vorhandenen Datensätze richtig klassifiziert sind. Dieses Problem kann 
man aber genauso schreiben als:

\begin{align*}
		\operatorname{minimiere}_{w,b} \quad & \norm{(w)} \\
		\operatorname{sodass} \quad &y_{i}(\langle w,x_{i} \rangle -b) \geq 1 \quad \forall i = 1,...,n ; \\
		&x \in \mathbb{R}^n; \quad y \in \{1; -1\}
		%\label{SVMProblemMin}
\end{align*}
%Genauso gemacht wie obendrüber, funktioniert aber nicht geaunso ??

Klassifikation:

\begin{align*}
	y_{i} =& \{ 1 \operatorname{falls} \langle w,x_{i} \rangle - b \geq 1\} \\
		& \{ -1 \operatorname{falls} \langle w,x_{i} \rangle \leq 1\}
	%\label{KlassifikationSVM}
\end{align*}
%????

Nachdem das Problem optimiert wurde, ist der Computer in der Lage, weitere Daten einzuordnen, vorrausgesetzt, es ist richtig optimiert.
Dennoch kann es bei dieser Art von SVMs zu einigen Problemen kommen.\\

\begin{itemize}
	\item Wenn die Datensätze nicht linear separierbar sind, d.h., es ist nicht möglich die beiden Datensätze mit einem linearen Element zu trennen\\
	
	\item Wenn die Daten nicht alle korrekt klassifiziert sind oder es Daten gibt, die in dem >>Bereich<< der anderen Seite liegen.
\end{itemize}

\subsubsection{Soft Margin SVMs}
Die Soft Margin SVMs können mit dem zweiten Problem umgehen.
Sie sind im Prinzip wie herkömmliche SVMs aufgebaut, nur dass sie eine gewisse Fehlertoleranz besitzen. Diese kommt durch eine Veränderung an der Formel zustande:

\begin{align*}
	\operatorname{minimiere}_{w,b} \quad & \norm{w} + C\sum_{i}{z_{i}} \quad \forall i = 1,...,n \\
	\operatorname{sodass} \quad & y_{i}(\langle w,x_{i} \rangle - b) \geq 1 \quad \forall i = 1,...,n; \\ 
	&x \in \mathbb{R}^n; y \in \{1; -1\} \quad z_{i} \geq 0 \\
	%\label{eq: SoftMarginSVMs}
\end{align*}

Dabei steht $z_{i}$ für die Grösse des Fehlers des Punktes $x_{i}$ und C ist eine Konstante, deren grösse bestimmt, wie stark die Fehler gewichtet werden, da die Konstante mit der Summe der Fehler multipliziert wird. Da die gesamte Gleichung aber minimiert werden soll sorgt ein hohes C dafür, dass das Problem schwerer optimiert werden kann. Die Konstante muss vom Mensch selbst gewählt werden, je nachdem, wie schwer Fehler gewichtet werden sollen.

\begin{dsafigure}
	\begin{center}
		\includegraphics[width=0.5\textwidth]{\media Figure_SVM.pdf}
		\caption{Links und rechts zur Trenngeraden befinden sich die parallelen Grenzen (gestrichelte Geraden). Ziel der Optimierung ist es, den Abstand zwischen den Grenzen zu maximieren, um den Normalenvektor $w$ zu bestimmen.}
		\label{FigSVM}
	\end{center}
\end{dsafigure}

\section{Neuronale Netze}
\author {Farhadiba Mohammed}

\subsection{Einführung}

Inspiriert von unserem Verständnis, wie das menschliche Gehirn lernt, benutzen Neuronale Netze Lernalgorithmen, welche besonders für praktische Anwendungen geeignet sind.
Dazu zählen Spracherkennung, Objekterkennung in Bildern und die Fähigkeit individuell passende Produkte vorzuschlagen, die dem Kunden gefallen könnten. 
Ein Neuronales Netz wird von mehreren Schichten aufgebaut. Ausgangsschicht ist dabei, die Datenschicht, auf die ein oder mehrere hidden layer folgen. Als Ausgabewert erhält man schließlich einen Vektor, welcher die Wahrscheinlichkeitsverteilung darstellt. Mit anderen Worten, wie wahrscheinlich es ist, dass das Ausgangsobjekt zu einer bestimmten Klasse gehört.

\begin{dsafigure}
	\begin{center}
		\includegraphics[width=0.5\textwidth]{\media Figure_NN.pdf}
		\caption{Ein vollständig verbundenes neuronales Netzwerk mit $i$ Eingängen und $k$ Ausgängen, bestehend aus $n$ Schichten mit jeweils $m$ >>Neuronen<<.}
		\label{FigNN}
	\end{center}
\end{dsafigure}

\subsection{Deep Learning}

Mit Deep Learning beschreibt man die Neuronalen Netze, die über mehr als einen >>versteckte Schicht<< verfügen. Damit ist gemeint, dass sich zwischen Ein- und Ausgabeschicht weitere Schichten befinden. 

Um die Anzahl an Parametern zu vergrößern, werden diese auch über nicht-lineare Funktionen miteinander verknüpft. Dadurch kann eine höhere Genauigkeit erzielt werden. Jedoch kann eine zu hohe Anzahl an Parametern auch dazu führen, dass das Netzwerk >>overfitted<<, d.h., zu sehr an den Trainingsdatensatz angepasst, wird. Dann kann das Netzwerk neue, fremde Daten nicht mehr korrekt klassifizieren. 

Wie wir in der Abbildung 2.3 sehen können ist ein typisches Fully connected Neuronales Netz abgebildet. Als Basis findet man unten die Datenschicht mit ihren Eingabewerten in Form eines Vektors $x_1 ... x_i$. Diese Werte werden nun mit einem jeweiligen Faktor $w$ multipliziert. Hierbei handelt es sich um ein Skalarprodukt. Das Ergebnis wird nun im ersten hidden layer abgespeichert. Darauf wird eine nicht lineare Funktion angewendet und das Ganze wird erneut mit einem Faktor $w$ multipliziert. Dieser Vorgang wiederholt sich so lange bis uns schließlich ein Vektor mit seinen Komponenten $p_1 .. p_1$ zurückgegeben wird, welcher als Wahrscheinlichkeitsverteilung interpretiert wird. Anhand eines konkreten Beispiels würde es folgendes bedeuten:
Nehmen wir an wir haben ein Bild und wollen ermitteln zu welcher Klasse Haus, Tisch oder Stuhl das daurauf abgebildete Objekt gehört. Unsere inputs wären demnach die Pixel des Bildes. Diese durchlaufen nun das Neuronale Netz und wir erhalten eine Wahrscheinlichkeitsverteilung als Rückgabewert, welche uns mitteilt, dass es am wahrscheinlichsten ist, dass das abgebildete Objekt ein Objekt der Klasse Stuhl ist. Demnach ist die Klassifizierung vollzogen.


\begin{align*}
a_j^{(k)} = f^{(k)} (\langle w_{ji}^{(k)}, a_i^{(k-1)}\rangle) = f^{(k)} (\sum_{i=1}^{m^{(k-1)}} w_{ji}^{(k)},a_i^{(k-1)} )
\end{align*}

\subsection{Convolutional Neural Networks}

\begin{dsafigure}
	\begin{center}
		\includegraphics[width=0.5\textwidth]{\media Figure_ConvNN}
		\caption{Ein Convolutional Neural Network (CNN) (dt.: >>faltendes neurales Netzwerk<<) mit vier Eingängen ($x_1, ..., x_4$) und zwei Ausgängen ($p_1, p_2$). Dazwischen befindet sich eine Schicht aus drei >>Neuronen<<, die als Filter wirkt.}
		\label{FigConvNN}
	\end{center}
\end{dsafigure}


Bei den Convolutional Networks handelt es sich um ein Neuronales Netz mit einer vereinfachten Struktur. Diese kommt dadurch zu Stande, dass nun kein Fully Connected Neuronales Netzwerk mehr vorliegt. Das bedeutet nicht jeder der Knoten ist mit jedem verbunden, sondern nur jeweils die lokalen drei.


\subsection{SVMs als Neuronale Netze}

Wir haben bereits SVMs kennengelernt. Im folgenden wird verdeutlicht, wie man diese auch als Neuronales Netz ausdrücken kann. Charakteristisch für dieses Neuronale Netz ist, dass es nur einen Layer besitzt:

\begin{align*}
\min\limits_{b,w,z} \norm{w} + C  \sum_{i=1}^{n} z
\end {align*}

sodass:

\begin{align*}
&y_i (<w,x> - b) \geq 1 - z_i ; z_i \geq 0 \\
\implies &z_i \geq 1- y_i (\langle w,x \rangle - b) \text{ und } z_i \geq 0
\end {align*}


\begin{align*}
\min &\norm{w} + C \cdot \sum_{i=1}^{n} \max (0, 1-y_i (\langle w,x_i\rangle-b) \\
 = &\norm{w} + C \sum_{i=1}^{n} \ell ( 1-y_i (\langle w,x_i\rangle)-b) \\
=&\norm{w} + C \sum_{i=1}^{n} \ell (a_i)
\end{align*}


\subsection{Lineare Regression}

Die lineare Regression kann auch in Form von Neuronalen Netzen ausgedrückt werden. 


\begin{align*}
\min\limits_{w,b} 0,5 \sum_{i=1}^{n} ((\langle w,x_i\rangle- b) -y_i)^2
a_i = \langle w,x_i\rangle -b-y_i
\end{align*}




%\input{Kapitel_LineareAlgebra}

%\input{Kapitel-2-Kapitelüberschrift}
%\input{Kapitel-3-Kapitelüberschrift}
%usw.

%\nocite{*} % alle referenzen anzeigen, sogar wenn sie nicht im text zitiert sind
\bibliography{lit}{}
\bibliographystyle{plain}
\end{document}
